\chapter{1977年贵州卷}

\section{解答题}

\begin{problem}
将下式化成最简形式：\(\sqrt[4]{a^{-2}b^3}\cdot\sqrt[3]{a^4b^{-2}}\div\sqrt[6]{a^5b^{-2}}\).
\end{problem}

\begin{solution}
在实数范围内，原式有意义且分母不为零时，因式中含有 \(a^{-2}\) 和 \(b^{-2}\)，知 \(a\neq0\)，\(b\neq0\). 四次根式要求 \(a^{-2}b^3\geqslant0\)，从而 \(b>0\)；六次根式要求 \(a^5b^{-2}\geqslant0\)，从而 \(a>0\). 因此可以使用正数幂的运算律，将各根式化为十二次根式，得
\[
\begin{aligned}
\text{原式}
&=\frac{\sqrt[12]{a^{-6}b^9}\cdot\sqrt[12]{a^{16}b^{-8}}}{\sqrt[12]{a^{10}b^{-4}}} \\
&=\sqrt[12]{\frac{a^{-6}b^9\cdot a^{16}b^{-8}}{a^{10}b^{-4}}}
=\sqrt[12]{\frac{a^{10}b}{a^{10}b^{-4}}}
=\sqrt[12]{b^5}.
\end{aligned}
\]
\end{solution}

\begin{problem}
将下式化成最简形式：\(\log_2\cos\frac{\pi}{4}-\log_2\sin\frac{\pi}{6}\).
\end{problem}

\begin{solution}
因为 \(\cos\frac{\pi}{4}=\frac{\sqrt2}{2}=2^{-1/2}>0\)，\(\sin\frac{\pi}{6}=\frac12=2^{-1}>0\)，所以两个对数都有意义. 代入对数的幂运算公式，得
\[
\text{原式}
=\log_2 2^{-1/2}-\log_2 2^{-1}
=-\frac12-(-1)=\frac12
\]
\end{solution}

\begin{problem}
已知 \(\tan\alpha=\dfrac{2ab}{a^2-b^2}\)，求 \(\sin\alpha\).
\end{problem}

\begin{solution}
由题意知 \(a^2-b^2\neq0\). 先由平方关系计算余弦的平方：
\[
\begin{aligned}
\cos^2\alpha
&=\frac{1}{1+\tan^2\alpha}
=\frac{1}{1+\dfrac{4a^2b^2}{(a^2-b^2)^2}} \\
&=\frac{(a^2-b^2)^2}{(a^2-b^2)^2+4a^2b^2}
=\frac{(a^2-b^2)^2}{(a^2+b^2)^2}.
\end{aligned}
\]
由于 \(a^2+b^2>0\)，可写成
\(\cos\alpha=\varepsilon\frac{a^2-b^2}{a^2+b^2}\)，\(\varepsilon\in\{1,-1\}\).
于是
\[
\sin\alpha=\tan\alpha\cos\alpha
=\frac{2ab}{a^2-b^2}\cdot\varepsilon\frac{a^2-b^2}{a^2+b^2}
=\varepsilon\frac{2ab}{a^2+b^2}
\]
因此，在未给出 \(\alpha\) 所在象限时，\(\sin\alpha=\pm\dfrac{2ab}{a^2+b^2}\)；若已知象限，则按象限确定正负号.
\end{solution}

\begin{problem}
已知 \(m+\dfrac{1}{m}=2\)，求

\begin{enumerate}
\item \(m^2+\dfrac{1}{m^2}=?\)；
\item \(m^3+\dfrac{1}{m^3}=?\).
\end{enumerate}
\end{problem}

\begin{solution}
因为分式有意义，所以 \(m\neq0\). 将已知等式乘以 \(m\)，得
\[
m^2-2m+1=0
\]
即 \((m-1)^2=0\)，所以 \(m=1\). 因而
\begin{enumerate}
\item \(m^2+\dfrac{1}{m^2}=1^2+\dfrac{1}{1^2}=2\)；
\item \(m^3+\dfrac{1}{m^3}=1^3+\dfrac{1}{1^3}=2\).
\end{enumerate}
\end{solution}

\begin{problem}
以直角三角形 \(ABC\) 的一直角边 \(AB\) 为直径作圆，此圆与斜边 \(AC\) 交于 \(D\)，过 \(D\) 引圆的切线交 \(BC\) 于 \(E\). 试证 \(EB=CE\).

\bitmapfigure[max width=5.8cm]{img/1977/guizhou/q5_fig1.png}
\end{problem}

\begin{solution}
连接 \(BD\). 因为 \(AB\) 是圆的直径且 \(D\) 在圆上，由圆周角定理，\(\angle ADB=90^\circ\). 又因 \(A\)，\(D\)，\(C\) 共线，所以 \(\angle CDB=90^\circ\).

由 \(\angle ABC=90^\circ\)，知 \(AB\perp BC\). 圆心在直径 \(AB\) 的中点上，故半径在 \(B\) 点与直线 \(BC\) 垂直，\(BE\) 是圆在 \(B\) 点的切线. 题设中的 \(DE\) 也是切线，因此从圆外一点 \(E\) 引圆的两条切线，切线段相等，即
\[
ED=EB
\]
于是三角形 \(EDB\) 为等腰三角形，设
\[
\angle EDB=\angle DBE=\delta
\]
再设 \(\angle CDE=\gamma\)，\(\angle DCE=\eta\). 由于 \(\angle CDB=90^\circ\)，有 \(\gamma+\delta=90^\circ\)；由于 \(C\)，\(E\)，\(B\) 共线，且三角形 \(CDB\) 在 \(D\) 点直角，有 \(\eta+\delta=90^\circ\). 因此 \(\gamma=\eta\)，三角形 \(CDE\) 的两角相等，从而
\[
CE=ED
\]
结合 \(ED=EB\)，得到 \(EB=CE\).
\end{solution}

\begin{problem}
求到点 \(A\left(1,2\right)\) 和点 \(B\left(4,1\right)\) 等距离的点的轨迹方程.
\end{problem}

\begin{solution}
设动点为 \(P(x,y)\). 由等距条件，
\[
PA=PB
\]
使用两点间距离公式，得
\[
\sqrt{(x-1)^2+(y-2)^2}=\sqrt{(x-4)^2+(y-1)^2}
\]
两边均为非负数，可以平方，得到
\[
(x-1)^2+(y-2)^2=(x-4)^2+(y-1)^2
\]
展开并整理：
\[
x^2-2x+1+y^2-4y+4=x^2-8x+16+y^2-2y+1
\]
所以 \(6x-2y-12=0\)，即
\[
3x-y-6=0
\]
因此所求轨迹是直线 \(3x-y-6=0\)，也可写成 \(y=3x-6\).
\end{solution}

\begin{problem}
用长 \(8\mathrm{~m}\) 的木料做一矩形窗框，问长和宽如何选择，才能使窗子透过的光线最多？
\end{problem}

\begin{solution}
设矩形窗子的宽为 \(x\mathrm{~m}\)，长为 \(y\mathrm{~m}\)，其中 \(x>0\)，\(y>0\). 木料长度等于窗框周长，所以
\(2(x+y)=8\)，\(y=4-x\)，\(0<x<4\).
窗子透过的光线多少与窗子的面积成正比，故只需使面积
\[
S=xy=x(4-x)=-(x-2)^2+4
\]
最大. 因为 \(-(x-2)^2\leqslant0\)，且当且仅当 \(x=2\) 时取等号，所以 \(S\) 的最大值为 \(4\mathrm{~m^2}\). 此时 \(y=4-x=2\).

因此长和宽都应取 \(2\mathrm{~m}\)，窗子为边长 \(2\mathrm{~m}\) 的正方形.
\end{solution}

\begin{problem}
求 \(\displaystyle\lim\limits_{n\to\infty}\sin a\left(1+\cos a+\cos^2a+\cdots+\cos^na\right)\)（其中 \(a\) 为任意实常数，\(0<a<\dfrac{\pi}{2}\)）.
\end{problem}

\begin{solution}
由 \(0<a<\dfrac{\pi}{2}\)，有 \(0<\cos a<1\). 对有限项等比数列求和，
\[
1+\cos a+\cos^2a+\cdots+\cos^na=\frac{1-\cos^{n+1}a}{1-\cos a}
\]
由于 \(0<\cos a<1\)，所以 \(\displaystyle\lim\limits_{n\to\infty}\cos^{n+1}a=0\). 因而
\[
\begin{aligned}
\text{原式}
&=\lim\limits_{n\to\infty}\frac{\sin a\left(1-\cos^{n+1}a\right)}{1-\cos a} \\
&=\frac{\sin a}{1-\cos a}
=\frac{1+\cos a}{\sin a}
=\cot\frac{a}{2}.
\end{aligned}
\]
\end{solution}

\begin{problem}
已知方程 \(x^2+px+q=0\) 的两个根为 \(\tan\alpha\) 与 \(\tan\beta\)，求 \(\sin^2\left(\alpha+\beta\right)+p\sin\left(\alpha+\beta\right)\cos\left(\alpha+\beta\right)+q\cos^2\left(\alpha+\beta\right)\) 的值（其中 \(p\)，\(q\) 为任意实常数）.
\end{problem}

\begin{solution}
令 \(u=\tan\alpha\)，\(v=\tan\beta\). 由根与系数的关系，
\(u+v=-p\)，\(uv=q\).
记 \(S=\alpha+\beta\). 先讨论 \(q\neq1\) 的情形. 因为 \(\tan\alpha\) 与 \(\tan\beta\) 有意义，\(\cos\alpha\cos\beta\neq0\)，且
\[
\cos S=\cos\alpha\cos\beta(1-uv)=\cos\alpha\cos\beta(1-q)\neq0
\]
故 \(\tan S\) 有意义，并且
\[
\tan S=\frac{u+v}{1-uv}=-\frac{p}{1-q}
\]
原式可写为
\[
\cos^2S\left(\tan^2S+p\tan S+q\right)
\]
又 \(\cos^2S=\dfrac{1}{1+\tan^2S}\)，所以
\[
\text{原式}=\frac{\tan^2S+p\tan S+q}{1+\tan^2S}
\]
代入 \(\tan S=-\dfrac{p}{1-q}\)，分子为
\[
\frac{p^2}{(1-q)^2}-\frac{p^2}{1-q}+q
=\frac{q\left(p^2+(1-q)^2\right)}{(1-q)^2}
\]
而分母为
\[
1+\frac{p^2}{(1-q)^2}=\frac{p^2+(1-q)^2}{(1-q)^2}
\]
由于 \(q\neq1\)，有 \(p^2+(1-q)^2>0\)，所以约去相同因子后，原式等于 \(q\).

当 \(q=1\) 时，\(uv=1\)，从而 \(\cos S=\cos\alpha\cos\beta(1-uv)=0\). 此时 \(\sin^2S=1\)，且含有 \(\cos S\) 的另外两项均为零，所以原式仍为 \(1=q\).

综上，所求值为 \(q\).
\end{solution}

\begin{problem}
设有直线 \(L\) 与曲线 \(c\)：\(L:y=ax-1\)，\(c:r\sin\theta=\sin2\theta\). 问 \(L\) 与 \(c\) 在什么条件下相交，相切或不相交？
\end{problem}

\begin{solution}
在非原点处，由极坐标与直角坐标的关系 \(x=r\cos\theta\)，\(y=r\sin\theta\)，\(r^2=x^2+y^2\)，原极坐标方程给出
\[
y=r\sin\theta=2\sin\theta\cos\theta=\frac{2xy}{x^2+y^2}
\]
原点也满足原极坐标方程（取 \(\theta=0\)）. 因此不能约去 \(y\)，否则会漏掉 \(y=0\) 这一整支. 两边相乘并整理，得
\[
y\left(x^2+y^2-2x\right)=0
\]
所以曲线 \(c\) 是两部分的并集：
\(c_1:y=0\)，\(c_2:x^2+y^2-2x=0,\qquad\text{即}\quad (x-1)^2+y^2=1\).

先考察直线 \(L\) 与 \(c_1\) 的关系. 当 \(a=0\) 时，\(L:y=-1\)，它与 \(c_1\) 不相交；当 \(a\neq0\) 时，令 \(y=0\)，得交点
\[
P\left(\frac1a,0\right)
\]

再考察直线 \(L\) 与圆 \(c_2\) 的关系. 代入圆方程，得
\[
(x-1)^2+(ax-1)^2=1
\]
即
\[
(1+a^2)x^2-2(a+1)x+1=0
\]
其判别式为
\[
\Delta=\left[-2(a+1)\right]^2-4(1+a^2)=8a
\]
因此，当 \(a>0\) 时，直线与圆有两个交点；当 \(a=0\) 时，直线与圆相切于 \((1,-1)\)；当 \(a<0\) 时，直线与圆没有交点.

合并两部分可得：
\begin{enumerate}
\item 当 \(a<0\) 时，\(L\) 与 \(c\) 在 \(\left(\frac1a,0\right)\) 相交，且不与圆支相交；
\item 当 \(a=0\) 时，\(L\) 与 \(c\) 相切于 \((1,-1)\)；
\item 当 \(a>0\) 时，\(L\) 与圆支有两个交点，并且还与 \(y=0\) 有交点 \(\left(\frac1a,0\right)\). 当 \(a=\frac12\) 时，该点是圆支交点 \((2,0)\) 之一，故共有两个不同交点；当 \(a>0\) 且 \(a\neq\frac12\) 时，共有三个不同交点.
\end{enumerate}
因此，按原极坐标方程的完整轨迹理解，\(a\neq0\) 时均有相交情形，\(a=0\) 时相切，没有不相交的情形. 约去 \(y\) 只能得到圆支，不能据此代替原曲线的完整分类.
\end{solution}

\section{参考题}

\begin{problem}
设某平板上任意一点的电势由 \(E=\ln\rho\) 给出（\(\rho\geqslant1\)），这里 \(\rho\) 是以原点为起点的矢径. 求在点 \(\rho=5\)，\(\theta=\arcsin0.8\) 处，沿此点至 \(\rho=10\)，\(\theta=\dfrac{\pi}{2}\) 的方向 \(E\) 的变化率.
\end{problem}

\begin{solution}
因为 \(\theta=\arcsin0.8\) 是第一象限角，所以 \(\sin\theta=0.8\)，\(\cos\theta=0.6\). 起点的直角坐标为
\[
P=(5\cos\theta,5\sin\theta)=(3,4)
\]
终点的直角坐标为
\[
Q=\left(10\cos\frac\pi2,10\sin\frac\pi2\right)=(0,10)
\]

\solutionfigure{\bitmapfigure[width=4.2cm]{img/1977/guizhou/q11_solution_fig1.png}}

把 \(E=\ln\rho\) 写成直角坐标函数，得
\[
E(x,y)=\ln\sqrt{x^2+y^2}
\]
因此
\(\frac{\partial E}{\partial x}=\frac{x}{x^2+y^2}\)，\(\frac{\partial E}{\partial y}=\frac{y}{x^2+y^2}\).
从 \(P\) 指向 \(Q\) 的方向向量为
\(\overrightarrow{PQ}=(-3,6)\)，\(\left|\overrightarrow{PQ}\right|=3\sqrt5\)
所以对应的单位方向向量为
\[
\bs{u}=\frac{\overrightarrow{PQ}}{\left|\overrightarrow{PQ}\right|}
=\left(-\frac1{\sqrt5},\frac2{\sqrt5}\right)
\]
于是，方向导数为
\[
\begin{aligned}
D_{\bs{u}}E(P)
&=\left(\frac{\partial E}{\partial x},\frac{\partial E}{\partial y}\right)_{(3,4)}\cdot\bs{u} \\
&=\left(\frac3{25},\frac4{25}\right)\cdot\left(-\frac1{\sqrt5},\frac2{\sqrt5}\right)
=\frac{\sqrt5}{25}.
\end{aligned}
\]
因此，所求变化率为 \(\dfrac{\sqrt5}{25}\).
\end{solution}

\begin{problem}
证明 \(S_1=1\)，\(S_2=1+\dfrac{1}{2^2}\)，\(S_3=1+\dfrac{1}{2^2}+\dfrac{1}{3^2}\)，\(\ldots\)，\(S_n=1+\dfrac{1}{2^2}+\dfrac{1}{3^2}+\cdots+\dfrac{1}{n^2}\)，\(\ldots\) 是收敛数列.
\end{problem}

\begin{solution}
对任意 \(n\geqslant1\)，有
\[
S_{n+1}-S_n=\frac{1}{(n+1)^2}>0
\]
所以数列 \(\{S_n\}\) 单调递增.

当 \(n\geqslant2\) 时，对每个 \(k=2\)，\(3\)，\(\ldots\)，\(n\)，都有 \(k^2>k(k-1)>0\)，从而
\[
\frac1{k^2}<\frac1{k(k-1)}=\frac1{k-1}-\frac1k
\]
因此
\[
\begin{aligned}
S_n
&=1+\sum_{k=2}^{n}\frac1{k^2}
<1+\sum_{k=2}^{n}\left(\frac1{k-1}-\frac1k\right) \\
&=1+\left(1-\frac12\right)+\left(\frac12-\frac13\right)+\cdots+\left(\frac1{n-1}-\frac1n\right) \\
&=2-\frac1n<2.
\end{aligned}
\]
而 \(S_1=1<2\)，所以整个数列有上界 \(2\). 单调递增且有上界的数列必收敛，故 \(\{S_n\}\) 是收敛数列.
\end{solution}
